\subsection*{Numerical and Computational Mathematics Notation}

\begingroup
\small

\noindent\emph{Status.} The numerical and computational chapters are mostly placeholders outside the numerical linear algebra already covered in the Linear and Operator Structures part. These entries reserve the intended notation without adding PDE or machine-learning conventions.

\notationgroup{Floating Point, Approximation, and Error}

\noindent\begin{tabularx}{\textwidth}{@{}p{0.38\textwidth}X@{}}
\notationentry{\operatorname{fl}(x),\ u,\ \varepsilon_{\mathrm{mach}}}{floating-point rounding, unit roundoff, and machine epsilon}
\notationentry{x_h,\ e_h,\ \mathcal{O}(h^p)}{discretized approximation, discretization error, and asymptotic order}
\notationentry{\widehat{x},\ \Delta x,\ r=b-A\widehat{x}}{computed approximation, perturbation or error in $x$, and residual}
\end{tabularx}

\notationgroup{Monte Carlo and Stochastic Computation}

\noindent\begin{tabularx}{\textwidth}{@{}p{0.38\textwidth}X@{}}
\notationentry{X_1,\ldots,X_N,\ \widehat{\mu}_N}{samples and Monte Carlo estimator}
\notationentry{\operatorname{Var}(\widehat{\mu}_N),\ \operatorname{SE}(\widehat{\mu}_N)}{variance and standard error of an estimator}
\notationentry{w_i,\ \widehat{I}_N}{importance weights and estimated integral or expectation, when introduced}
\end{tabularx}

\endgroup
